\section{Background}
\label{sec:background}
This section reviews the technical foundations relevant to policy-private consent enforcement on blockchains. We first introduce the cryptographic access-control mechanisms commonly used in blockchain-based healthcare data sharing. We then review zero-knowledge proofs and the Data Use Ontology (DUO)~\cite{dyke2016consent}, which together provide the basis for the policy representation and verification mechanisms used in our approach.

\subsection{Cryptographic Access Control in Blockchain Healthcare}
Fine-grained consent management in blockchain-based healthcare data sharing is commonly built on cryptographic access-control primitives, chiefly Ciphertext-Policy Attribute-Based Encryption (CP-ABE) and Proxy Re-Encryption (PRE)~\cite{bethencourt2007ciphertext, ateniese2006improved,ullah2023survey}.
In CP-ABE, the encrypting party embeds an attribute-based access policy into the ciphertext, so that only requesters whose credentials satisfy the policy can decrypt the protected data~\cite{ibraimi2009ciphertext}.
PRE lets a semi-trusted intermediary transform a ciphertext encrypted for one party into one decryptable by a designated requester, delegating access without exposing the plaintext~\cite{ateniese2006improved}.
Both primitives remain the building blocks of current blockchain healthcare designs: recent schemes combine PRE with on-chain records to delegate access to medical records and IoMT data~\cite{pei2024proxy,zhou2024secure}, and CP-ABE variants with hidden policies and attributes, distributed key generation, or multi-authority hierarchies address the exposure of the access structure itself~\cite{yang2024hamed,ding2026badsabe,wu2023blockchain}.
We cite the original constructions because the later systems inherit their trust and disclosure properties, which is what matters for this paper: these primitives establish who may obtain the data, and how they relate to \emph{policy} confidentiality on a public chain is discussed in Section~\ref{sec:related-work}.

\subsection{Zero-Knowledge Proofs}
Zero-knowledge proofs (ZKPs) are cryptographic protocols that allow a prover to convince a verifier that a statement is true without revealing any information beyond the validity of the statement~\cite{groth2016size}.
Modern ZKP systems support verifiable computation over private inputs, so that complex predicates can be proven without disclosing the underlying data.
In blockchain settings, ZKPs are widely used to reconcile public verifiability with privacy: privacy-focused systems such as Zerocash and Tornado Cash apply them to conceal transaction details or provide transaction anonymity~\cite{sasson2014zerocash, nadler2023tornado}.
These properties make ZKPs a natural building block for proving that a requester satisfies a hidden consent policy, which is the direction our approach develops.

\subsection{The Data Use Ontology}
\label{subsec:duo}
The Data Use Ontology (DUO) provides a standardized vocabulary for expressing biomedical data-sharing consent conditions, and its terms are widely adopted in genomics and biobanking to enable semantic interoperability~\cite{lawson2021data, dyke2016consent}.
A DUO-based policy combines permissions and modifiers---for example, general research use (GR), disease-specific research (DS), or a requirement that ethics-committee approval be obtained (IRB)---into the condition under which access is granted.

\paragraph{DUO Term Categories}
Although DUO offers a rich set of permissions and modifiers, its terms are heterogeneous in how they must be evaluated for compliance.
To structure their handling in a zero-knowledge setting, we group DUO terms into four categories:

\begin{itemize}
	\item \textbf{Boolean-only:} simple yes/no conditions that can be represented directly without external validation, e.g.\ ``no restriction'' (NRES), general research use (GR), or ``health/medical/biomedical research'' (HMB).

	\item \textbf{Ontology-linked:} terms whose evaluation involves consistency checks against external vocabularies or classifications, e.g.\ disease-specific research (DS), population origins or ancestry research (POA), or clinical care use (CC).

	\item \textbf{Obligation-required:} terms that impose procedural or temporal obligations, e.g.\ ``publication required'' (PUB), ``publication moratorium'' (MOR), ``time limit on use'' (TS), or ``return to database'' (RTN).

	\item \textbf{Credential-required:} terms whose evaluation relies on attested proof of requester status or approvals, e.g.\ ``collaboration required'' (COL), ``ethics approval required'' (IRB), ``project-specific restriction'' (PS), or ``institution-specific restriction'' (IS).
\end{itemize}

\noindent
Table~\ref{tab:duo-categories} summarizes representative DUO terms under this classification. Some terms reduce to simple Boolean flags, while others rely on ontology references, external obligations, or credential attestations. Later sections describe how these heterogeneous requirements are harmonized into a unified privacy-preserving verification process.

\begin{table}[H]
	\centering
	\caption{Representative DUO terms grouped by evaluation category.}
	\label{tab:duo-categories}
	\begin{tabular}{@{}>{\raggedright\arraybackslash}p{0.20\linewidth}>{\raggedright\arraybackslash}p{0.30\linewidth}>{\raggedright\arraybackslash}p{0.42\linewidth}@{}}
		\toprule
		\textbf{Category}                                                         & \textbf{Example Terms} & \textbf{Description} \\
		\midrule
		Boolean-only                                                              &
		NRES (DUO:0000004), GR (DUO:0000042), HMB (DUO:0000006)                   &
		Direct yes/no conditions, no external validation required.                                                                \\
		\addlinespace
		Ontology-linked                                                           &
		DS (DUO:0000007), POA (DUO:0000011), CC (DUO:0000043)                     &
		Evaluated against external ontologies such as ICD-10 or MONDO.                                                            \\
		\addlinespace
		Obligation-required                                                       &
		PUB (DUO:0000019), MOR (DUO:0000024), TS (DUO:0000025), RTN (DUO:0000029) &
		Impose obligations such as publication duties, time limits, or data return.                                               \\
		\addlinespace
		Credential-required                                                       &
		COL (DUO:0000020), IRB (DUO:0000021), PS (DUO:0000027), IS (DUO:0000028)  &
		Depend on credentials, approvals, or institutional/project-specific restrictions.                                         \\
		\bottomrule
	\end{tabular}
\end{table}
